\section*{\emoji{woman-teacher} Excel-Projekt}

Im Rahmen dieses Auftrags werden Sie ein eigenes Excel-Projekt erstellen. Das Thema des Projekts können Sie sich frei auswählen. Alternativ finden Sie untenan eine Liste mit Ideen für Ihr Projekt. Ihr Endresultat sollte aus folgenden Teilen bestehen:

\begin{itemize}
	\item \textit{Eine} Excel-Datei (\texttt{.xlsx}), in welcher all Ihre Daten, Resultate, Grafiken, und weiteres Material gespeichert sind.
	\item Ein Tab ``\texttt{Daten}'' pro Datenset, das Sie anlegen (z.B. Liste von Personen, von Waren, etc.)
	\item Ein Analyse-Teil, entweder in den Daten-Tabs oder in separaten Analyse-Tabs. Hier sollten Sie mindestens einmal (in inhaltlich sinnvoller) Weise die Anwendung von Sortier- und Filter-Funktionen aufzeigen.
	\item Tabs ``Resultate'', die Ihre Resultate enthalten. Diese Tabs sollten mindestens 3 Pivot-Tabellen enthalten, die spannende Resultate aufzeigen, sowie 2 mindestens Grafiken unterschiedlicher Art.
\end{itemize}

Zudem verfassen Sie zu Ihrem Projekt ein maximal 2 Seiten langes Begleitdokument, mit welchem eine aussenstehende Person Ihr Projekt von A bis Z reproduzieren kann. Das  Dokument sollte folgende Aspekte beinhalten:
\begin{itemize}
	\item Datenquellen: Welche Daten verwenden Sie und weshalb haben Sie sich für genau diese Daten entschieden? Falls Ihre Daten nicht vollständig oder fehlerhaft sind, welche Einschränkungen ergeben sich daraus?
	\item Daten-Bearbeitungen: Mussten Sie die Daten aussortieren, aufräumen oder korrigieren für Ihre Analysen?
	\item Analysen: Wie haben Sie Ihre Daten analyisiert, um Schlüsse / Einsichten daraus zu generieren? Wie sind Sie vorgegangen (Methode, noch keine Resultate)
	\item Resultate: Was haben Sie herausgefunden (Grafiken hier einfügen, allenfalls Pivot Tables zeigen).
	\item Offene Punkte: Welche Aspekte würden Sie gerne noch weiter ausbauen, sofern Sie mehr Zeit für das Projekt erhalten hätten?
\end{itemize}

\subsection*{Bewertung}

Ihr Projekt wird nach folgenden Kriterien bewertet:
\begin{itemize}
	\item (2 Punkte) \textbf{Excel-Kriterien}: Wurden alle zuoberst genannten Kriterien erfüllt (mehrere Daten-Quellen, Analysen, Pivot-Tabellen etc.)? Sind die unterschiedlichen Teile klar in separate Excel-Tabs aufgeteilt?
	\item (2 Punkte) Sind die gesammelten oder erfundenen Daten von hoher Qualität und Ausführlichkeit (viele Kolonnen und Zeilen)? Können daraus vielfältige, sinnvolle Einsichten gewonnen werden?
	\item (2 Punkte) \textbf{Resultate}: Sind Ihre Resultate anschaulich und klar präsentiert? Haben Sie sowohl mehrere Pivot-Tabellen wie auch mehrere Grafiken sinnvoll verwendet?
	\item (4 Punkte) \textbf{Bericht}: Klar und grammatikalisch fehlerfrei verfasst, enthält die erwähnten Aspekte in sinnvoller Struktur. Der Bericht erlaubt es, Ihr Vorgehen zu verstehen, so dass Ihr Projekt von einer aussenstehenden Person nachgebaut werden kann.
\end{itemize}

\newpage
\subsection*{Mögliche Themen (Beispiele)}

Idealerweise finden Sie ein Thema, das Sie persönlich interessiert, und mit welchem Sie einen für Sie interessanten Lebensbereich abdecken können. Falls Ihnen Ideen für ein Excel-Projekt fehlen, sind hier einige Vorschläge. Die Daten zu den jeweiligen Vorschläge dürfen Sie erfinden oder Daten aus dem Internet verwenden, sofern Ihre Daten-Quellen klar angegeben sind.

\subsubsection*{Organisation Klassen-Woche}
Sie sind eine Lehrperson und organisieren mit einer weiteren Lehrperson gemeinsam ein Klassenlager organisieren. Eine Ihrer Hauptaufgaben besteht aus der Menuplanung für die gesamte Woche. Zwecks Budget- und Einkaufsplanung schlagen Sie vor, die Menuplanung in Excel zu machen. Ihr Vorgehen hierzu ist wie folgt:

\begin{itemize}
	\item Sammeln der relavanten Daten (1. Teil): Daten zur Klasse: Wie viele Personen sind in der Klasse, welche Essens-Einschränkungen gibt es (Allergien, Vegetarisch / Vegan)?
	\item Sammeln von Menu-Ideen: Sammeln Sie einige Ideen von Gerichten, welche Sie gerne kochen möchten. Erstellen Sie jeweils Varianten für die jeweiligen Essens-Allergien / -Einschränkungen. Gliedern Sie jedes Gericht in die Lebensmittel (Zutaten) auf, welche für das Gericht benötigt werden. Berechnen Sie pro Zutat die benötigte Menge. Beschränken Sie sich in Ihrer Analyse auf die ersten beiden Tage. 
	\item Sammeln der relevanten Daten (2. Teil): Sammeln Sie online, z.B. auf \url{https://coop.ch}, Daten zu den Lebensmitteln die Sie für Ihre Gerichte verwenden wollen. Tragen Sie folgende Daten in eine neue Tabelle ein: Lebensmittel-Name, Menge (Gewicht) pro Verkaufseinheit, Preis pro Einheit, Web-Adresse (URL).
	\item Berechnen Sie nun den Preis jedes Gerichts, indem Sie die eingeplante Menge pro Zutat und Gericht mit dem Preis multiplizieren. 
	\item Geben Sie in einer Pivot-Tabelle aus, welche Gerichte wie viel kosten, und zeigen Sie für jedes Gericht auf, welche Zutat am teuersten ist (pro Person).
\end{itemize}
\textbf{Bonus}: Berechnen Sie für jedes Ihrer Gerichte auch die Nährwerte (Kalorien, Protein, Kohlenhydrate, Fett).

\subsubsection*{Fitness-Studio}
Sie besitzen ein eigenes Fitness-Studio. Um sicherzustellen, dass Ihr Fitness-Studio auch in Zukunft rentabel bleibt, möchten Sie jeden Monat möglichst viele neue Abonnements abschliessen. Erfassen Sie folgendes:

\begin{itemize}
	\item Mitgliedschafts-Arten in Ihrem Fitnessstudio (z.B. Monatlich, Jährlich), sowie den Preis der jeweiligen Mitgliedschaften
	\item Mögliche, gebührenpflichtige Extras Ihres Fitnessstudios (mind. 5 Angebote): Z.B: Abonnements auf Gruppen-Fitness-Kurse, Getränke, Sport-Snacks etc.	
	\item Mitglieder Ihres Studios: Name, Alter, Geschlecht, Wohnort, Mitglieds-Art, Dauer der Mitgliedschaft (Abschluss-Datum der ersten Mitgliedschaft), Ablauf der aktuellen Mitgliedschaft (im Datums-Format!)
	\item Zeigen Sie in einem Tab diejenigen Mitglieder an, deren Mitgliedschaft in weniger als 2 Monaten abläuft. Diesen Mitgliedern sollen personalisierte Verlängerungs-Angebote per Mail geschickt werden.
\end{itemize}

Sie entscheiden sich, ein Spezial-Extra-Angebot zu erstellen, welches denjenigen Mitgliedern, die sich für den Kauf des Extra-Angebots entscheiden, 50 Gramm Proteinpulver pro Tag bereitstellt. Um profitabel zu sein, suchen Sie sich die günstigste Quelle für Proteinpulver aus. Erfassen Sie dafür in einer Datenbank verschiedene Proteinpulver-Produkte von Anbietern wie z.B. Coop, Migros und Lee-Sports und vergleichen Sie deren Preis \textit{pro Gramm}. Welches Pulver ist am günstigsten pro Gramm? Wie viel muss Ihr Abo kosten, falls jedes Mitglied im Schnitt 3 mal pro Woche ein Produkt bezieht und Sie mindestens 20\% Profit machen wollen? Stellen Sie Ihre Resultate als Pivot-Tabellen und als Excel-Grafiken dar.

\textbf{Bonus}: Sammeln (bzw, erfinden) Sie Daten zu den Fitness-Studios in Ihrer Umgebung und vergleichen Sie deren Erfolg (z.B. Anzahl Mitglieder, Einkommen pro Mitglied etc.) mit Ihrer eigenen Firma. Schlagen Sie basierend auf den Daten eine Strategie vor, um Mitglieder der anderen Fitnessstudios abzuwerben.


\subsubsection*{Weitere Daten-Quellen}
\begin{itemize}
	\item \url{https://opendata.swiss/de/dataset}
	\item \url{https://www.kaggle.com/datasets}
	\item \url{https://data.world/datasets/open-data}
\end{itemize}